\begin{frame}{Circuito 2: efeito da carga resistiva}
  \small
  \begin{columns}[T,onlytextwidth]
    \begin{column}{0.47\textwidth}
      \begin{block}{Equação nodal e função de transferência}
        \[
          \frac{V_i-V_o}{R}=sCV_o+\frac{V_o}{R_L}.
        \]
        \[
          H_2(s)=\frac{R_L}{R+R_L+sCRR_L}
          =\frac{K_2}{1+s\tau_2}.
        \]
        \[
          K_2=\frac{R_L}{R+R_L},\qquad
          \tau_2=(R\parallel R_L)C.
        \]
      \end{block}
    \end{column}
    \begin{column}{0.49\textwidth}
      \begin{alertblock}{Projeto nominal com carga de $22\,\mathrm{k}\Omega$}
        \[
          R\parallel R_L=13{,}2\,\mathrm{k}\Omega.
        \]
        \[
          K_2=0{,}4\quad(-7{,}96\,\mathrm{dB}).
        \]
        \[
          f_{c2}=\frac{1}{2\pi(R\parallel R_L)C}
          \approx120{,}57\,\mathrm{Hz}.
        \]
        O ganho diminui, mas a frequência de corte aumenta:
        o capacitor passa a enxergar uma resistência equivalente menor.
      \end{alertblock}
    \end{column}
  \end{columns}
\end{frame}
